My weld-seam HMI joined modified-DH kinematics, analytical forward and inverse kinematics, an OpenGL robot model, and perception software, while my later CR5 work showed how command timing can destabilize an otherwise coherent robot loop. These projects gave me useful integration experience but revealed a capability gap: I need a more rigorous framework for connecting control, industrial automation, and manufacturable system design. CUHK's MSc in Mechanical and Automation Engineering offers that combination. If available in my intake, MAEG 5725 Control and Industrial Automation and MAEG 5755 Robotics would deepen the control and automation reasoning behind the interfaces I have implemented.

I am also drawn to MAEG 5160 Design for Additive Manufacturing because intelligent machines should be evaluated not only as software and control stacks, but also as physical systems whose geometry and production constraints shape integration. Subject to supervision and availability, MAEG 5910 M.Sc. Project would let me bring these perspectives into one bounded prototype, tracing a design from mechanical representation through controller behavior and system validation. I would contribute experience locating mismatches among coordinate conventions, software boundaries, and execution timing, while preparing to build robotics and intelligent-manufacturing systems whose mechanical design, automation logic, and verification reinforce one another.
